% AI-produced by GPT-6 Astra (OpenAI, ChatGPT Pro mode). No priority claim.
% This fragment follows A.tex and uses its lemmas, notation, and claim environment.
\subsection{A realization theorem and the unrestricted localizer classification}

\begin{claim}{A.13}{Proved}{Realization with control of every action}\label{A:general-realization}
Let $K$ be a nonempty connected small Kan complex. Define
\[
 \mathcal A(K)=\Delta_{\mathrm{inj}}/K.
\]
Thus an object is a finite nonempty linearly ordered set $S$ with a simplicial map $\Delta^S\to K$, and an arrow is an increasing injection preserving that map. Composition is composition of injections, and every simplex, including every degenerate simplex, is an object. Then
\[
 B\mathcal A(K)\simeq |K|.
\]
For every action on $\mathcal A(K)$, each connected component of its category of elements is either weakly contractible or weakly equivalent, over $|K|$, to a connected covering space of $|K|$. Conversely every connected covering space of $|K|$ occurs as an action component. In particular, if $K$ is simply connected, every action component is weakly contractible or its projection to $\mathcal A(K)$ is a weak equivalence. For $n\geq2$,
\[
 \mathcal A(K)\in\mathcal Q_n
 \quad\Longleftrightarrow\quad
 |K|\text{ is }(n-1)\text{-connected}.
\]
\end{claim}
\begin{proof}
The comma contraction in Claim~\ref{A:cocycle-category} applies to every simplicial set: prepending a vertex labelled zero contracts the category of increasing-domain factorizations of a simplex. It gives
$B(\Delta_{\mathrm{inj}}/K)\simeq B(\Delta/K)\simeq|K|$, naturally in $K$. We prove the stronger all-action assertion.

\emph{Functorial amalgamation over nonempty bases.}
Choose, once and for all, one filler for every horn map into $K$. There are only set many choices. For a span $P\to X,P\to Z$ of finite $K$-labelled simplices with $P\ne\varnothing$, take the tagged set union $U=X\amalg_P Z$, ordered by placing the $X$-side before the $Z$-side in each gap of $P$. Put $p=\min P$. The simplices wholly in one side retain their existing labels. We define the labels of the mixed simplices in increasing dimension as follows.

For every mixed nonempty subset $T\subset U\setminus\{p\}$, assume the simplices on $\{p\}\cup(T\setminus\{t\})$ have been defined for every $t\in T$. Their labels fit together to a horn on the ordered set $\{p\}\cup T$, missing the face $T$. Apply the chosen filler for that horn. Use its full simplex as the label on $\{p\}\cup T$ and its missing face as the label on $T$. The compatibility of the horn follows inductively from the face identities. The label on the missing face has the already prescribed labels on all its proper faces, since those are the corresponding intersections with the other horn faces. This defines compatible labels on every subset of $U$, hence one simplex $\Delta^U\to K$ extending both inputs.

An embedding $X\to X'$ under $P$, together with $\id_Z$, preserves the order, the anchor $p$, and the distinction between mixed and one-sided subsets. Induction on the size of $T$ shows that it encounters exactly the same horn map and hence exactly the same chosen filler. Thus this construction is a functor of $X$, with natural embeddings $X\to U\leftarrow Z$. Every nonempty full cosieve of $P/\mathcal A(K)$ is consequently contractible by the argument of Claim~\ref{A:cocycle-amalgam}. This is a specified functor, not a noncanonical extension assertion.

We also use a strict extension fact. If $L\subset\Delta^U$ is a finite subcomplex and a map $L\to K$ is nullhomotopic, it extends literally to $\Delta^U$. Indeed, the restriction
\[
 \operatorname{Map}(\Delta^U,K)\longrightarrow\operatorname{Map}(L,K)
\]
is a Kan fibration. Its image on vertices is a union of connected components: lift a finite edge path from an extendible constant map to the given map. Nullhomotopy places the given map in the component of a constant map, which extends. The mapping-space assertion is precisely \cite[Proposition I.5.2 and Corollary I.5.3(2)]{GoerssJardine1999}, whose proof turns the lifting square into the pushout-product of a horn inclusion with $L\hookrightarrow\Delta^U$. In particular, a map on a union of simplices with contractible components extends when $K$ is connected; and a map on a union homotopy equivalent to a circle extends when $K$ is simply connected. In the disconnected-domain case choose paths between its finitely many component constants first.

\emph{Countable simply connected $K$.}
The finite labelled simplices form a countable class, with hereditary substructures, joint embeddings, and strong amalgamation. The first follows by restricting labels. Joint embeddings follow from the preceding extension fact on a disjoint union of two simplices, and amalgamation follows either from the explicit construction or from extension on two simplices meeting in a simplex. Construct a countable rich ultrahomogeneous limit $M$ as in Claim~\ref{A:cocycle-supports}. The language has a relation for every simplex label and the linear order; a finite structure is determined by its top-dimensional label.

The proof of that claim now gives, verbatim,
\[
 \operatorname{Fix}_M(G_B)=B,\qquad
 \langle G_B,G_C\rangle=G_{B\cap C},
 \qquad G=\operatorname{Aut}(M).
\]
Here is the extension step that replaces the cocycle calculation. For two disjoint-off-$D$ neighboring copies $B_0,B_1$, prescribe the three labelled simplices on $B_0\cup B_1$, $B_0\cup C'$, and $B_1\cup C'$. Their union has contractible nerve if $D\ne\varnothing$ and circle nerve if $D=\varnothing$. Its map to $K$ is nullhomotopic, so the strict extension fact applies. For the finite chain of intermediate copies, prescribe the existing \emph{whole} simplex on $B_0\cup B_1$ and the separate new copies; their intersections with that existing piece and with one another are exactly $D$. The union has contractible components. All order-gap interpolation and finite coset-path arguments are unchanged. This verifies the hypotheses at the precise two places where Claim~\ref{A:cocycle-supports} used cohomology.

Thus every open subgroup of $G$ is a finite pointwise stabilizer. Applying the explicit quotient-action and completion arguments of Claim~\ref{A:all-actions}, using the functorial amalgam just constructed, proves that every action component is contractible or projects by a weak equivalence to $\mathcal A(K)$. This proves the simply connected countable case without assuming any result about elimination of imaginaries.

\emph{Countable connected $K$ and its universal cover.}
Put $\Gamma=\pi_1|K|$, and let $\widetilde K\to K$ be the simplicial universal covering. Both $\Gamma$ and $\widetilde K$ are countable. The deck action of $\Gamma$ is free on the vertices of $\widetilde K$. Construct the rich limit $M$ for $\widetilde K$ and write $G=\operatorname{Aut}(M)$ as above.

Let $\widehat G$ be the group of increasing permutations $g$ of $M$ for which there is a deck transformation $\gamma\in\Gamma$ such that the label on $gS$ is $\gamma$ times the label on $S$, for every finite nonempty $S$. The degree $\gamma$ is unique, since the deck action is free on vertices. Every degree occurs: twisting all labels of $M$ by $\gamma$ gives another countable rich limit of the same class, and back-and-forth gives the required semilinear isomorphism. We obtain
\[
 1\longrightarrow G\longrightarrow\widehat G
   \xrightarrow{\rho}\Gamma\longrightarrow1.
\]
Give $\widehat G$ the pointwise convergence topology. The subgroup fixing any one vertex lies in $G$, so $G$ is open and $\Gamma$ is discrete. Every finite semilinear embedding extends to $\widehat G$: first use a permutation of the desired degree, and then extend the remaining ordinary partial isomorphism by homogeneity.

The category $\widehat A$ of finite nonempty subsets of $M$, with semilinear increasing embeddings, is equivalent to $\mathcal A(K)$. The functor forgets the chosen lift of each label. It is essentially surjective because every simplex of $K$ has a lift that occurs in $M$. It is fully faithful because two lifts of one simplex differ by a unique deck transformation; therefore every increasing injection preserving labels downstairs has a unique semilinear degree upstairs.

For a nonempty finite $B$, its pointwise stabilizer in $\widehat G$ is exactly $G_B$. If $U\leq\widehat G$ is open, the simply connected case gives $U\cap G=G_F$ for a finite $F$, possibly empty. Since $G$ is normal, conjugation by $u\in U$ gives $G_{uF}=G_F$. The fixed-point identity forces $uF=F$. If $F\ne\varnothing$, preservation of the finite linear order makes $u$ fix every point of $F$. Its degree is then trivial because the deck action on vertex labels is free. Hence $U=G_F$. If $F=\varnothing$, then $U$ contains $G$ and has the form $\rho^{-1}(H)$ for a subgroup $H\leq\Gamma$. This exhausts all open subgroups, including those arising as imaginary stabilizers.

Repeat the two explicit homotopy-final action comparisons of Claim~\ref{A:all-actions} with $\widehat G$. They require only the directed poset of finite subsets of $M$, extension of finite embeddings, and contractibility of nonempty cosieves in slices, all proved above. For an orbit $\widehat G/G_F$ with $F\ne\varnothing$, the fixed-coset calculation identifies the resulting action with $\widehat A(F,-)$. Its category of elements is contractible. An orbit $\widehat G/\rho^{-1}(H)$ is the set $\Gamma/H$; $G_B$ fixes it entirely, and an arrow of degree $\gamma$ acts by left multiplication by $\gamma$. This is a locally constant action, and its category of elements is the corresponding covering of $B\widehat A$.

To identify its monodromy precisely, first take $H=1$. The element category of the regular $\Gamma$-action is equivalent to $\mathcal A(\widetilde K)$: an object $(B,\gamma)$ is sent to the labelled simplex $\gamma^{-1}\sigma_B$. An arrow $f$ of degree $\delta$, with $\delta\gamma=\gamma'$, preserves these transformed labels because
$\gamma'^{-1}\delta\sigma_B=\gamma^{-1}\sigma_B$. This functor is fully faithful and essentially surjective by the lifting description. Its classifying space is connected and simply connected. Therefore this regular covering is universal and the monodromy is the specified group $\Gamma$, up to the harmless base-point choice. The orbit $\Gamma/H$ consequently gives the connected cover with subgroup $H$. This proves the countable connected case, including control of the projection map.

\emph{Arbitrary small connected $K$.}
The poset of countable connected Kan subcomplexes $K_i\subset K$ is filtered and covers $K$. Given countably many simplices, close under faces and degeneracies, connect their vertices to one chosen vertex by countably many finite paths, and choose fillers for all horns in successive countable stages. This produces such a subcomplex. The same procedure starting from a union proves directedness.

Restrict an arbitrary action $L$ to each $\mathcal A(K_i)$. Fix an object $(a,l)$ of a component $D$ of $\El L$, and restrict to the cofinal directed subposet of $i$ containing $a$. Let $D_i$ be the component containing $(a,l)$. For $i\leq j$, its inclusion maps into $D_j$. These are inclusions of nerves, and
\[
 ND=\bigcup_i ND_i,
 \qquad N\mathcal A(K)=\bigcup_iN\mathcal A(K_i).
\]
The first equality holds because each zigzag connecting an object to $(a,l)$ uses finitely many simplex labels. Homotopy groups commute with these filtered unions of simplicial sets: a finite simplicial representative and any finite simplicial nullhomotopy are present at a single later stage, after the usual finite subdivisions or fibrant-replacement construction.

Each projection $BD_i\to B\mathcal A(K_i)$ is injective on $\pi_1$, because its source is contractible or is a connected cover by the countable case. The same is true in the union: a loop whose image is nullhomotopic is represented together with that nullhomotopy at a later stage, where injectivity kills it. If $BD$ is weakly contractible there is nothing further to prove. Otherwise choose a nonzero class in some $\pi_m(BD)$, $m\geq1$, and a stage representing it. At every later stage that class remains nonzero, since it survives in the union. Thus every later $D_i$ is of covering type. Its projection is an isomorphism on all homotopy groups of degree at least two. Filtered passage gives the same assertion for $BD\to|K|$. Together with injectivity on $\pi_1$, the covering-space lifting criterion and Whitehead's weak-equivalence criterion identify $BD$, over $|K|$, with the connected cover corresponding to its fundamental-group image. This argument does not assume that components cannot merge: their merging is incorporated in the specified components $D_i$.

Finally, any simplicial connected cover $K_H\to K$ induces a discrete opfibration $\mathcal A(K_H)\to\mathcal A(K)$: a lift of one face uniquely extends over its whole simplex. Its associated covariant set action has category of elements $\mathcal A(K_H)$, whose classifying space is the desired cover. Thus every connected cover occurs. The assertion about $\mathcal Q_n$ follows: necessity uses the singleton action, and sufficiency uses simple connectivity of an $(n-1)$-connected space and the all-action description.
\end{proof}

\begin{claim}{A.14}{Proved}{Exactly which arbitrary localizers give an intermediate condition}\label{A:global-intermediate}
Let $W\supseteq W_\infty$ be any fundamental localizer. Define a class of ordinary connected homotopy types by
\[
 \mathscr H_W=
 \{X:\ \Sigma X_H\text{ is }W\text{-aspherical for every connected cover }
 X_H\to X\}.
\]
This definition involves ordinary covering spaces and suspensions, with no quantification over categorical actions. Write $\mathscr C_k$ for the nonempty $k$-connected homotopy types. For every $n\geq2$, the suspension condition is strictly between two consecutive tower levels precisely as follows:
\[
 \mathcal Q_{n+1}\subsetneqq\mathcal S_W\subsetneqq\mathcal Q_n
 \quad\Longleftrightarrow\quad
 \mathscr C_n\subsetneqq\mathscr H_W\subsetneqq\mathscr C_{n-1}.
 \tag{A.14.1}
\]
Equivalently, the necessary and sufficient four conditions on $W$ are:
\begin{enumerate}
\item Every $n$-connected $X$ has $W$-aspherical suspension.
\item If all connected covers of $X$ have $W$-aspherical suspensions, then $X$ is $(n-1)$-connected.
\item Some $(n-1)$-connected but not $n$-connected $X$ has $W$-aspherical suspension.
\item Some $(n-1)$-connected $Y$ has a suspension that is not $W$-aspherical.
\end{enumerate}
Whenever these conditions hold, $\mathcal A(\operatorname{Sing}X)$ and $\mathcal A(\operatorname{Sing}Y)$ are completely specified witnesses for the two strict inclusions. In the adjacent interval $W_{n+1}\subseteq W\subseteq W_n$, this reduces to the concrete group criterion
\[
 \mathcal R_W\text{ contains a nonzero abelian group and omits a nonzero abelian group},
\]
with no countability restriction. If $\mathcal R_W=\{0\}$, then $\mathcal S_W=\mathcal Q_{n+1}$; if $\mathcal R_W=\Ab$, then $\mathcal S_W=\mathcal Q_n$.
\end{claim}
\begin{proof}
We first record why covers introduce no extra categorical assumption. For a connected category $D$, a connected cover of $BD$ is realized by a locally constant covariant set action of $D$, obtained from its fundamental groupoid. Its category of elements has the required covering homotopy type. If $D$ is already an action component of $A$, the composite-action construction of Claim~\ref{A:cone} realizes this as another action component of $A$. Hence
\[
 A\in\mathcal S_W
 \quad\Longleftrightarrow\quad
 BD\in\mathscr H_W\text{ for every action component }D\text{ of }A.
 \tag{A.14.2}
\]
The reverse direction uses the identity cover.

If $\mathscr C_n\subseteq\mathscr H_W$, (A.14.2) gives $\mathcal Q_{n+1}\subseteq\mathcal S_W$. Conversely, apply that categorical inclusion to $\mathcal A(\operatorname{Sing}X)$ for any $n$-connected $X$ and use Claim~\ref{A:general-realization}. Its base is an action component, so $X\in\mathscr H_W$. Similarly, $\mathscr H_W\subseteq\mathscr C_{n-1}$ implies $\mathcal S_W\subseteq\mathcal Q_n$ by (A.14.2); the converse again follows from the realization theorem, since $X\in\mathscr H_W$ makes every action component of $\mathcal A(\operatorname{Sing}X)$ have $W$-aspherical suspension.

These arguments also preserve strictness. A space in either set-theoretic difference gives a categorical witness by the realization theorem. A category in either categorical difference supplies an action component in the corresponding spatial difference, using (A.14.2). This proves (A.14.1). Since $n\geq2$, all spaces in $\mathscr C_{n-1}$ are simply connected and have only their identity connected cover. This gives the displayed four-condition formulation and its witnesses.

In the adjacent interval, formula (A.12.1) already gives the two inclusions. By Claim~\ref{A:general-realization}, applied to $K(\Lambda,n)$ for an arbitrary small abelian group $\Lambda$, the countable restriction in the witnesses of Claim~\ref{A:intermediate-localizers} can be removed. Thus a nonzero member of $\mathcal R_W$ gives strictness at the lower end and a nonzero nonmember gives strictness at the upper end. Necessity follows directly from (A.12.1): a category in the lower difference has some nonzero $H_n(BD)$ in $\mathcal R_W$; a category in the upper difference has some $H_n(BD)$ outside $\mathcal R_W$, which is necessarily nonzero because $K(0,n+1)$ is contractible. The two equality statements follow from the same formula and Hurewicz. This completes the classification for arbitrary $W$ as a condition on its ordinary suspension-null homotopy types, while the adjacent-interval case has the simpler abelian-group form.
\end{proof}
